\subsection{Submanifolds in symplectic manifolds}

In these notes a submanifold means an {\it embedded\/} submanifold.
By the local $k$-slice a condition, a subspace $Q\subseteq M$ is an embedded $k$-manifold iff there is a cover of $Q$ by open subsets of $M$
such that $Q$'s image in these charts is an open subset of $\bR^k\subseteq\bR^n$.

\bdefn

    Let $(M,\omega)$ be a symplectic manifold.
    A submanifold $Q\subseteq M$ is {\emphcolor symplectic}, {\emphcolor isotropic}, {\emphcolor coisotropic}, or {\emphcolor Lagrangian}
    if for every $q\in Q$, $T_qQ\subseteq T_qM$ is a symplectic/isotropic/coisotropic/Lagrangian subspace.

\edefn

It is clear then that $Q$ is Lagrangian iff it is isotropic and of dimension $\frac12\dim M$.

\bexam

    Let $Q\subseteq L$ be an arbitrary submanifold.
    Then define its {\emphcolor annihilator} (or {\emphcolor conormal})
    $$ TQ^\perp = \set{(q,\phi)\in T^*L}[q\in Q,\ \phi|_{T_qQ}=0] . $$
    This is a vector bundle over $Q$, and each fiber has dimension $\dim L-\dim Q$ (the dimension of $\ann(T_qQ)\subseteq\ann(T_qL)$).
    Thus $TQ^\perp$'s rank is $\dim L-\dim Q$, so its dimension is $\dim L$.

    Notice that $\lambda_\can$ vanishes on $TQ^\perp$: $\lambda_\can(q,\phi)=\phi\circ\d\pi_{(q,\phi)}$, and $\d\pi_{(q,\phi)}$ lands in $T_qQ$ which
    is annihilated by $\phi$.
    Thus $Q$ is isotropic and its dimension is half of $T^*L$'s, so it is Lagrangian.

    Taking $Q=L$, we get the zero section
    $$ TL^\perp = \set{(q,0)\in T^*L}[q\in L] , $$
    and taking $q\in L$ and $Q=\set q$, we have $T_qQ=0$ and thus
    $$ TL^\perp = T_q^*L . $$
    So both the zero section and the fibers of $T^*L$ are Lagrangian submanifolds.

\eexam

\bexam

    Let $(M,\omega)$ be a symplectic manifold, then its product $M\times M$ has a natural symplectic structure by $(-\omega)\oplus\omega$.
    Then $M\times*,*\times M$ are symplectic clearly.
    The diagonal is Lagrangian, as it is isotropic and of half of $M\times M$'s dimension.

\eexam

\bexam

    All $1$-dimensional submanifolds are isotropic, as all $1$-dimensional subspaces are isotropic.
    This means that all subspaces of codimension $1$ are coisotropic, and thus all hypersurfaces (i.e.\ submanifolds of codimension $1$) are coisotropic.

\eexam

\bprop[name={closedformlagrange}]

    Let $\sigma\in\Omega^1(M)$ be a $1$-form.
    Then its graph $\Gamma_\sigma\subseteq T^*M$ is Lagrangian iff $\sigma$ is closed.

\eprop

\bproof

    The graph $\Gamma_\sigma$ is the image of the embedding $\sigma\colon M\to T^*M$.
    Thus $-\d\lambda_\can$ vanishes on $\Gamma_\sigma$ iff
    $$ 0 = \sigma^*(-\d\lambda_\can) = \d(\sigma^*\lambda_\can) = \d\sigma . $$
    The last equality is proven as follows.
    Write $\sigma$ locally by $\sigma=\sum_i\sigma_i\d x_i$, so $\sigma\colon M\to T^*M$ is the map $(x_1,\dots,x_n)\to(x_1,\dots,x_n,\sigma_1,\dots,\sigma_n)$.
    The pullback of $\lambda_\can$ along $\sigma$ then is
    $$ \sigma^*\lambda_\can(\partial_i) = \lambda_\can\parens{\pdv{x_i} + \sum_j \pdvof{\sigma_j}{x_i}\pdv{\sigma_j}} , $$
    and as $\lambda_\can=\sum_i\sigma_i\d x_i$, this is equal to $\sigma_i$.
    Thus $\sigma^*\lambda=\sigma$ as desired.
    \qed

\eproof

\bprop[name={symplagrange}]

    Let $(M,\omega)$ be symplectic and $\psi\colon M\to M$ a diffeomorphism.
    Then $\psi$ is a symplectomorphism iff its graph
    $$ \Gamma_\psi = \set{(q,\psi q)}[q\in M] \subseteq M\times M $$
    is a Lagrangian submanifold of $(M,\times M,(-\omega)\oplus\omega)$.

\eprop

\bproof

    By taking a coordinate chart $U\subseteq M$, then $U\times U$ is a coordinate chart of $M\times M$.
    We see then that $(U\times U)\cap\Gamma_\psi$ is the fiber of $0$ under the map $U\times U\to U$ given by $(q,p)\mapsto \psi(q)-p$.
    This map has full rank and so we see that this defines an embedded manifold structure on $\Gamma_\psi$ of dimension $\dim U=\dim M$.
    Furthermore the tangent space of $\Gamma_\psi$ at a point of $U\times U$ is the kernel of this map, which is
    $$ T_{(q,\psi q)}\Gamma_\psi = \set{(v,d\psi v)}[v\in T_qM] = \Gamma_{\d_q\psi} . $$

    Now $\psi\colon(W,\omega)\to(W,\omega)$ is a symplectomorphism iff $\Gamma_\psi\subseteq W\times W$ is Lagrangian.
    Indeed, it is Lagrangian iff it is isotopic iff $\omega(\psi w)-\omega(v)=0$ for all $w\in W$, i.e.\ iff $\psi$ is a symplectomorphism.

    Applying this to the above computation $\psi$ is a symplectomorphism if $\d_q\psi$ is for all $q\in M$ iff $\Gamma_{\d_q\psi}$ is Lagrangian
    for all $q\in Q$ iff $\Gamma_\psi\subseteq M\times M$ is Lagrangian.
    \qed

\eproof

If $F\inj E$ is a subbundle, then we can form the quotient bundle $E/F$ whose fibers are $E_p/F_p$.
If $Q\subseteq M$ is a submanifold, the quotient bundle $TM|_Q/TQ$ is called the {\emphcolor normal bundle} $\nu_Q$.
When $M$ is Riemannian, this is isomorphic to $TQ^\perp$ via the canonical isomorphism $W^\perp\cong V/W$ for finite-dimensional inner product spaces.

If $M$ is symplectic and $Q$ a symplectic submanifold then as in the Riemannian case $\nu_Q$ may be identified with $TQ^\omega$ and thus carries
a symplectic structure.
In the following theorem $\nu_Q$ refers to $TQ^\omega$.

\bthrm

    Let $(M_j,\omega_j)$ be symplectic manifolds for $j=0,1$ and $Q_j\subseteq M_j$ be compact symplectic submanifolds.
    Furthermore suppose that there is an isomorphism
    $$ \Phi\colon\nu_{Q_0} \to \nu_{Q_1} $$
    of the symplectic normal bundles which covers a symplectomorphism $\phi\colon(Q_0,\omega_0)\to(Q_1,\omega_1)$.
    Then $\phi$ extends to a symplectomorphism $\psi\colon(U_0,\omega_0)\to(U_1,\omega_1)$ for open neighborhoods $Q_i\subseteq U_i$
    such that $\d\psi$ restricts to $\Phi$ on $\nu_{Q_0}$.

\ethrm

\bproof

    By the tubular neighborhood theorem there are convex neighborhoods of the zero section of $\nu_{Q_i}$ which map
    diffeomorphically to a neighborhood $U_i$ of $Q_i$.
    Let $\exp_i$ be these diffeomorphisms and define
    $$ \phi' = \exp_1\circ\Phi\circ\exp_0^{-1} . $$
    On $\nu_{Q_0}$ we have that $\d\phi'=\Phi$ because on the zero section $\exp_i$'s differential is the identity.
    Now consider $\omega_1'=(\phi')^*\omega_1$, so that on $TM|_Q$ we have
    $$ \omega_1' = \Phi^*\omega_1 = \omega_0 . $$

    So we can instead consider $\omega_0,\omega_1'$.
    These are two symplectic forms on $U_0$ which agree on the compact submanifold $Q$.
    The result then follows from \refmath{moserisotopy}.
    \qed

\eproof

\bnote

    This result can be summarized as ``the symplectic structure around a compact symplectic submanifold is determined by the symplectic structure of its
    symplectic normal bundle.''

\enote

\blemm

    Let $\phi\colon M\to N$ be smooth between manifolds of the same dimension.
    Suppose that there is a compact submanifold $Q\subseteq M$ such that $\phi|_Q$ is injective and for every $q\in Q$, $\phi$ has full rank.
    Then there is a neighborhood $Q\subseteq U$ such that $\phi\colon U\to N$ is an embedding.

\elemm

\bproof

    Because being of full rank is local, for every point $q\in Q$ there is a neighborhood on which $\phi$ has full rank.
    Potentially shrinking this neighborhood, $\phi$ is a diffeomorphism there.
    It is sufficient for $\phi$ then to be injective in a neighborhood of $Q$, as injective local diffeomorphisms are embeddings.
    Because smooth manifolds are metrizable (e.g.\ by a Riemannian distance function), the following claim proves the lemma.

    Let $Q\subseteq M$ be a compact subset of a metric space, and $\phi\colon M\to N$ a continuous function.
    Suppose that around every point of $Q$ there is a neighborhood on which $\phi$ is injective, and that $\phi$ is injective on $Q$.
    Then there is a neighborhood of $Q$ on which $\phi$ is injective.

    Define
    $$ \epsilon = \inf\set{d(x,Q)}[\hbox{There is a $y$ such that $\phi x=\phi y$ and $d(y,Q)<d(x,Q)$}] . $$
    This $\epsilon$ is nonzero, otherwise there are sequences $x_n\neq y_n$ such that $d(x_n,Q),d(y_n,Q)\to0$ and $\phi(x_n)=\phi(y_n)$.
    Because $Q$ is compact there are $a_n,b_n\in Q$ such that $d(x_n,Q)=d(x_n,a_n)$ and $d(y_n,Q)=d(y_n,b_n)$.
    And again because $Q$ is compact there are subsequences of $a_n,b_n$ which converge, so we can assume that the sequences as a whole converge
    $a_n\to a$ and $b_n\to b$.
    Then $x_n\to a$ and $y_n\to b$ as well, so $\phi(a)=\phi(b)$, meaning $a=b$ as $\phi$ is injective on $Q$.
    But this means eventually $x_n,y_n$ lie in the open neighborhood of $a$ for which $\phi$ is injective, so $x_n=y_n$ a contradiction.

    So $\epsilon>0$, and $\phi$ is injective on the open tube
    $$ U = \set{x\in M}[d(x,Q)<\epsilon] . \qed $$

\eproof

\bthrm

    Let $L\subseteq M$ be a compact Lagrangian submanifold.
    Then there exists a neighborhood $U_0\subseteq T^*L$ of the zero section, a neighborhood $L\subseteq V$,
    and a symplectomorphism $\phi\colon(U_0,\omega_\can)\to(V,\omega)$ which restricts to the identity on $L$ (i.e.\ it maps the zero section to $L$).

    Here $\omega_\can=-\d\lambda_\can$ is the canonical symplectic structure on $T^*L$ (see \refmath[remark]{cotangentsymp}).

\ethrm

\bproof

    Let $J\colon TM\to TM$ be a complex structure compatible with $\omega$ (which exists by \refmath[proposition]{homotopysympbund}).
    Denote the associated Riemannian metric by $g_J(v,u)=\omega(v,Ju)$.
    For $q\in Q$, $J_qT_qQ\subseteq T_qM$ is the orhogonal complement of $T_qQ$ with respect to $g_J$ and is Lagrangian.
    In general if $\Lambda\subseteq V$ is a Lagrangian subspace, then $J\Lambda$ and $\Lambda$ have zero intersection as $\omega(v,Jv)>0$ for $v\in\Lambda$,
    and $g_J(Jv,u)=\omega(Jv,Ju)=\omega(v,u)=0$ for $v,u\in\Lambda$.

    Define
    $$ \Phi_q \colon T^*_qL \to T_qL $$
    to be the isomorphism induced by $g_J$, meaning
    $$ g_J(\Phi_q\phi,\bul) = \phi(\bul) . $$
    And consider $\phi\colon T^*L\to M$ by
    $$ \phi(q,\phi) = \exp_q(J_q\Phi_q(\phi)) , $$
    where $\exp\colon\E\subseteq TL\to M$ is the Riemannian exponential map associated to $g_J$.
    For $(v,\phi)\in T_{(q,0)}T^*L=T_qL\oplus T_q^*L$ we have
    $$ \d\phi_{(q,0)}(v,\phi) = v + J_q\Phi_q(\phi) . $$
    And thus if $v=(v_0,\phi_0)$ and $w=(v_1,\phi_1)$ are tangent vectors in $T_{(q,0)}T^*L$ then because $T_qL$ and $J_qT_qL$ are Lagrangian
    $$ \multline{
        \phi^*\omega_{(q,0)}(v,w) = \omega_q(v_0+J_q\Phi_q(\phi_0),v_1+J_q\Phi_q(\phi_1)) = \omega_q(v_0,J_q\Phi_q(\phi_1)) - \omega_q(v_1,J_q\Phi_q(\phi_0))\cr
        = g_J(v_0,\Phi_q(\phi_1)) - g_J(v_1,\Phi_q(\phi_0)) = \phi_1(v_0) - \phi_0(v_1) = \omega_\can(v,w) .
    } $$
    Thus the $2$-form $\phi^*\omega$ agrees with the symplectic structure $\omega_\can$ on the zero section of $T^*L$.
    Furthermore $\phi^*\omega$ is closed as the pullback of a closed form.

    So by \refmath{moserisotopy} there is a diffeomorphism $\psi\colon U_0\to U_1$ of neighborhoods of the zero section such that $\psi^*\phi^*\omega=\omega_\can$ and $\psi|_L=\id$.
    Furthermore because $\phi^*\omega$ is nondegenerate on the zero section, $\d\phi$ must have zero kernel on the zero section so that at points on the zero section $\phi$ has full rank.
    Because $T^*L$ and $M$ have the same dimension, by the above lemma $\phi$ is a diffeomorphism between from a neighborhood of the zero section to an open subset of $M$.
    Thus $\phi\circ\psi$ is the desired diffeomorphism.
    \qed

\eproof

\bnote

    This result shows that the symplectic structure around a compact Lagrangian submanifold is entirely determined by the diffeomorphism class of $L$.

\enote

\TODO{Whitney ($C^k$) topologies}

\bprop

    Let $(M,\omega)$ be symplectic, then a neighborhood of the identity in $\Symp(M)$ can be identified with a neighborhood of $0$ in the vector space of closed
    $1$-forms on $M$.

\eprop

\bproof

    Consider the Lagrangian submanifold $\Delta$ of the product $(M\times M,(-\omega)\oplus\omega)$.
    There is a canonical diffeomorphism $\Delta\cong M$ and thus an isomorphism $T^*\Delta\cong T^*M$.
    By the above theorem there is a neighborhood $\N(\Delta)\subseteq M\times M$ of $\Delta$, a convex neighborhood $\N(M_0)\subseteq T^*M$ of the zero section
    $M_0$, and a diffeomorphism $\Psi\colon\N(\Delta)\to\N(M_0)$ such that
    $$ \Psi(q,q) = (q,0),\qquad \Psi^*(\d\lambda_\can) = (-\omega)\oplus(\omega) $$
    (the first condition is $\Psi|_{M_0}=\id$, under the identification $\Delta\cong M$).

    Now suppose $\psi\in\Symp(M,\omega)$ is sufficiently $C^1$-close to the identity.
    By \refmath[proposition]{symplagrange} $\psi$'s graph is a Lagrangian submanifold of $\N(\Delta)$.
    Thus $L=\Psi(\Gamma_\psi)$ is a Lagrangian submanifold of $T^*M$.
    Furthermore $L$ is the image of $M\to T^*M$, $q\mapsto\Psi(q,\psi(q))$, which is $C^1$-close to the zero section.
    Thus $L$ is the graph of a $1$-form, and it is closed by \refmath[proposition]{closedformlagrange}.
    \qed

\eproof

